% ================================================================ chapter 6
\chapter{Conclusion and Future Work}
\label{ch:conclusion}

\section{Conclusion}
This thesis presented a unified forensic framework for the detection of
multimodal synthetic media that moves beyond binary classification to
fine-grained, per-modality attribution: determining independently whether the
video and the voice of a clip are real or AI-generated, assigning the clip to one
of four authenticity quadrants, and localizing manipulated segments in time. The
framework combines DAVID-Net, a disentangled cross-modal transformer, with
Quadrant-Aware Contrastive Pretraining, a generator-free pretraining strategy
that constructs all four quadrants from pristine data alone.
\todo{Summarize the headline empirical findings once available.} The complete
system is released as reproducible open-source code with a deployed public
detection service.

\section{Limitations}
\begin{itemize}
  \item The framework targets single-subject talking-head content; multi-person
        scenes and non-speech audio are out of scope.
  \item Detection is probabilistic; outputs are decision support, not proof.
  \item \todo{Add empirical limitations discovered during experiments.}
\end{itemize}

\section{Future Work}
\begin{itemize}
  \item Extension to fully generated (text-to-video) content without a real
        source subject.
  \item Continual adaptation to newly emerging generators using the QACP
        construction as a self-refreshing training signal.
  \item Multi-person scenes, overlapping speech, and non-frontal faces.
  \item Provenance integration (e.g., C2PA content credentials) alongside
        statistical detection.
  \item On-device and streaming (real-time) variants of the detector.
\end{itemize}
